% context: teacher solutions for 1-2CW_Velocity
% topic: average speed, unit conversion between ft/s, mi/h, m/s, km/h
\documentclass[12pt, twoside]{article}
\input{../preamble}
\usepackage{float}

\title{Physics}
\author{Chris Huson}
\date{September 2026}

\fancyhead[LE]{\thepage}
\fancyhead[RO]{\thepage}
\fancyhead[LO]{La Scuola d'Italia / Huson / Physics \\* 17 September 2026}

\begin{document}

\subsubsection*{1.2 Classwork Solutions: Speed, Velocity, and Unit Conversion}

\begin{enumerate}[itemsep=0.35cm]

\item A student walks 24.0 m in 18.5 s.
  \begin{enumerate}[itemsep=0.25cm]
    \item $v = \dfrac{d}{t} = \dfrac{24.0~\text{m}}{18.5~\text{s}} = 1.2972\ldots = \mathbf{1.30~\textbf{m/s}}$
    \item \textbf{Three significant figures.} Both measurements have three, and a quotient
      carries no more significant figures than the least precise measurement that went into it.
      Writing 1.29729 would claim a precision the stopwatch and the tape measure do not have.
  \end{enumerate}

\item Convert each speed.
  \begin{enumerate}[itemsep=0.35cm]
    \item $30.0~\dfrac{\text{km}}{\text{h}} \times \dfrac{1000~\text{m}}{1~\text{km}}
      \times \dfrac{1~\text{h}}{3600~\text{s}}
      = \dfrac{30{,}000}{3600}~\dfrac{\text{m}}{\text{s}} = 8.333\ldots
      = \mathbf{8.33~\textbf{m/s}}$

      \emph{The km cancel and the h cancel, leaving m/s. A useful benchmark: dividing
      km/h by 3.6 gives m/s.}
    \item $25.0~\dfrac{\text{m}}{\text{s}} \times \dfrac{1~\text{km}}{1000~\text{m}}
      \times \dfrac{3600~\text{s}}{1~\text{h}} = \mathbf{90.0~\textbf{km/h}}$
  \end{enumerate}

\item A fastball at 95 mi/h.
  \begin{enumerate}[itemsep=0.35cm]
    \item $95~\dfrac{\text{mi}}{\text{h}} \times \dfrac{5280~\text{ft}}{1~\text{mi}}
      \times \dfrac{1~\text{h}}{3600~\text{s}}
      = \dfrac{501{,}600}{3600}~\dfrac{\text{ft}}{\text{s}} = 139.33\ldots
      = \mathbf{139~\textbf{ft/s}}$
    \item $139.33~\dfrac{\text{ft}}{\text{s}} \times \dfrac{1~\text{m}}{3.28~\text{ft}}
      = \mathbf{42.5~\textbf{m/s}}$

      \emph{Either route works. Going through kilometers:
      $95 \times 1.61 = 153~\text{km/h}$, and $153 \div 3.6 = 42.5~\text{m/s}$.
      Note the full-precision 139.33 was carried into part (b) --- rounding to 139 first
      still gives 42.4, close but not identical. Round only at the end.}
  \end{enumerate}

\item Put all three in m/s before comparing.
  \begin{itemize}[itemsep=0.2cm]
    \item Car: $60~\dfrac{\text{mi}}{\text{h}} \times \dfrac{5280~\text{ft}}{1~\text{mi}}
      \times \dfrac{1~\text{h}}{3600~\text{s}} = 88.0~\text{ft/s}$, and
      $88.0 \div 3.28 = \mathbf{26.8~\textbf{m/s}}$
    \item Train: $95~\text{km/h} \div 3.6 = 26.39\ldots = \mathbf{26.4~\textbf{m/s}}$
    \item Motorcycle: $\mathbf{27~\textbf{m/s}}$ (already in m/s)
  \end{itemize}
  \vspace{0.2cm}
  Slowest to fastest: \textbf{train (26.4 m/s) $<$ car (26.8 m/s) $<$ motorcycle (27 m/s)}.

  \emph{The three are within 3\% of each other --- the point of the problem is that you
  cannot rank speeds until they are all in the same unit.}

\newpage

\item Sound travels 343 m/s; the delay is 4.0 s.

  $d = vt = (343~\text{m/s})(4.0~\text{s}) = 1372~\text{m} = \mathbf{1.4 \times 10^{3}~\textbf{m}}$
  \quad (2 s.f., because 4.0 s has two)

  $1372~\text{m} \times \dfrac{3.28~\text{ft}}{1~\text{m}} \times \dfrac{1~\text{mi}}{5280~\text{ft}}
  = \mathbf{0.85~\textbf{mi}}$

  \emph{This is the ``count the seconds'' rule of thumb: roughly 1 km every 3 seconds,
  or about 1 mile every 5 seconds.}

\item Climbing 5 floors, 18 steps each, 8.0 in per step.
  \begin{enumerate}[itemsep=0.3cm]
    \item $5~\text{floors} \times \dfrac{18~\text{steps}}{1~\text{floor}} = \mathbf{90~\textbf{steps}}$
    \item $90~\text{steps} \times \dfrac{8.0~\text{in}}{1~\text{step}} = 720~\text{in}$

      $720~\text{in} \times \dfrac{1~\text{ft}}{12~\text{in}} = \mathbf{60~\textbf{ft}}$
      \qquad
      $60~\text{ft} \times \dfrac{1~\text{m}}{3.28~\text{ft}} = 18.29\ldots = \mathbf{18~\textbf{m}}$

      \emph{The step counts are exact, so the 8.0 in riser sets the precision: two
      significant figures.}
    \item $t = \dfrac{90~\text{steps}}{1.2~\text{steps/s}} = \mathbf{75~\textbf{s}} = \mathbf{1~\textbf{min}~15~\textbf{s}}$
    \item $v = \dfrac{60~\text{ft}}{75~\text{s}} = 0.80~\dfrac{\text{ft}}{\text{s}}
      \times \dfrac{60~\text{s}}{1~\text{min}} = \mathbf{48~\textbf{ft/min}}$

      $v = \dfrac{18.3~\text{m}}{75~\text{s}} = \mathbf{0.24~\textbf{m/s}}$

      \emph{Check: $0.80 \div 3.28 = 0.244$ m/s --- the two routes agree. This is the
      \emph{vertical} speed; the distance actually walked along the stairs is longer,
      since each step also moves you forward.}
  \end{enumerate}

\item \textbf{Challenge.} The conversion itself is right:

  $30~\dfrac{\text{mi}}{\text{h}} \times \dfrac{1.61~\text{km}}{1~\text{mi}} = 48.3~\text{km/h}$

  What is wrong is the number of digits. The speedometer reads 30 mi/h --- two
  significant figures --- so the converted speed can carry at most two:
  $\mathbf{48~\textbf{km/h}}$. The seven digits in 48.28032 came out of the calculator,
  not out of the measurement; they claim the car's speed is known to a hundred-thousandth
  of a km/h. (The conversion factor is a rounded number too: 1 mi = 1.609344 km.)

  \emph{Converting a measurement never makes it more precise --- it only restates the
  same precision in a different unit.}

\end{enumerate}
\end{document}
